\section{Comunicare}
La comunicazione è la condivisione e scambio di informazioni. Si attua coi seguenti elementi fondamentali (più opzionalmente un terzo elemento che fornisce il mezzo di trasporto): \footnotesize\[[Codificatore/Trasmettitore] \longleftrightarrow [Trasporto] \longleftrightarrow [Ricevitore/Decodificatore]\] \normalsize Oggi le comunicazioni non devono avere limiti di spazio e tempo (approssimativamente) e essere standardizzate.

I mezzi trasmissivi ottimali sono la luce e i sengali elettromagnetici sono la soluzione ottimale.
\subsection{Terminologia}
\paragraph{Il Canale} è un'entità logica o fisica (diverse per tecnologia e principi fisici) che permette il trasporto dei singoli flussi informativi fra nodi remoti nello spazio.

Classificabili in Monodirezionali, Bidirezionali, Bidirezionali assimmetrici (quando un verso è usato più dell'altro), Punto-punto, Punto-multipunto Broadcast (tutti) o Multicast (sottoinsieme).

Un canale è definito dalla quantità di informazioni trasportabile e qualità del trasporto.
\subparagraph{Multicanale} è un unico collegamento che implementa più canali.
\paragraph{La Rete} è un insieme di canali condivisi per comunicazioni a richiesta con lo scopo del riuso dei canali. Una rete si compone di:
\begin{itemize}
	\item \textbf{Terminali}, sono l'interfaccia dell'utente alla rete, codificano e decodificano l'informazione in modo tale che sia comprensibile all'utente o alla rete.
	\item \textbf{Collegamenti}, permettono il trasferimento di uno o più flussi di informazioni fra punti remoti nello spazio.
	\item \textbf{Nodi di commutazione}, utilizzano i mezzi trasmissivi al fine di creare canali di comunicazione sulla base delle richieste degli utenti in modo automatico secondo regole prefissate.
\end{itemize}
\subparagraph{Rete gerarchica} organizza si più livelli la rete:
\begin{itemize}
	\item \textbf{Rete accesso}, collega i terminali a nodi periferici connessi a loro volta a nodi intermedi. (Topologia a stella).
	\item \textbf{Rete di trasporto}, collega a lunga distanza i nodi di transito. (Maglia completa).
\end{itemize}
Più si va nel kernel della rete più i canali si fanno "capienti" (perchè supportano un numero maggiore di utenti) e robusti rispetto alle periverie.
\paragraph{Topologie di rete} danno una descrizione geometrica di una rete tramite grafi (nodi e rami). Alcuni grafi:
\begin{itemize}
	\item \textbf{Maglia completa}, ha un collegamento per ogni coppia di nodi terminali ($N(N-1)/2$ collegamenti).

	Tanta resilienza e costo/complessità. Usato nel Kernel di una rete.
	\item \textbf{Stella}, ogni nodo ha un solo collegamento al nodo di comunicazione centrale attivo/passivo ($N$ collegamenti).

	Poca resilienza e costo/complessità. Usato nelle periferie di una rete.
	\item \textbf{Bus attivo/passivo}, Ogni nodo collegato a unico collegamento condiviso bidirezionale.
	\item \textbf{Anello}, ogni nodo collegato con soli altri due nodi monodirezionalmente o bidirezionalmente.
\end{itemize}
\subsection{Informazione}
Le categorie di sorgenti d'informazione sono voce e suoni, immagini e i dati. Queste categorie sono determinante dalla sorgente dell'informazione.
\paragraph{Il Segnale} è il risultato della trasformazione dell'informazione originale, prodotta dalla sorgente e appartenente a una delle tre classi, in una grandezza elettrica.
\subparagraph{La sorgente} può essere analogica (infiniti valori) o digitale (numero finito di valori, 0 e 1 tipicamente).

\paragraph{Segnale sinusoidale} è un segnale analogico molto semplice \[A\sin(2\pi ft + \phi)\] dove $A$ è l'ampiezza ossia l'intensità del segnale, $f$ è la frequenza ossia il tono del segnale (unità di misura l'Hz ossia il numero di oscillazioni al secondo) e $\phi$ è la fase (di poco interesse per noi).

\textbf{Spettogramma}: Rappresenta le sinusioni in un grafico di x=$f$ e  y=$A$) come segmenti verticali. Si usa la trasformata di Fourier per fare lo spettro di ogni segnale, potendoli rappresentare come somme di sinusoidi.
\begin{paracol}{2}
	\subparagraph{La Larghezza di banda} sono le frequnze comprese tra la $f_m$ e la $f_M$ nello spettro che un segnale assume. Si calcola facendo $f_M-f_m=B$.

	Più è grande la larghezza di banda più il segnale sarà complesso.
	\switchcolumn
	\subparagraph{La Banda passante} sono le frequenze trasmissibili sul mezzo trasmissivo e con la tecnologia scelta.
\end{paracol}

\medskip

\noindent La banda passante può essere inferiore alla larghezza di banda del segnale, e quindi filtrarne una parte, fintanto si comunichino informazioni utili e non uguali alle originali.
\subsection{Multi-informazione}
Abbiamo l'obbiettivo di trasportare sullo stesso mezzo trasmissivo più informazioni in canali differenti, ossia avere una comunicazione \textbf{intercanale}.
\paragraph{L'Interferenza} è il risultato distruttivo dalla somma irriversibile dei segnali (sofrapposizione degli spettri) se vengono trasmessi sullo stesso canale.
\paragraph{Il Multiplex} sposta le frequenze dall'intervallo originario a un altro e viceversa, in modo tale da non avere interferenze (ogni segnale non può sommarsi con altri avendo un intervallo proprio). Si usano le formule di prostaferesi. \[\cos(2\pi f_1t)\cos(2\pi f_2t)=\frac{1}{2}[\cos(2\pi f_ht)+\cos(2\pi f_lt)]\]\[f_h=f_1+f_2, f_l=f_1-f_2\]
Ora una rete manda un multiplex e nel mezzo trasmissivo abbiamo vari sotto canali.
\begin{itemize}
	\item \textbf{Modulazione AM}, modifica l'ampiezza dei segnali con somme e moltiplicazioni seguendo le formule proposte.
	\item \textbf{Modulazione FM}, modifica la frequenza dei segnali in modo più coplesso della AM ma con risultati migliori e per questo ancora usata.
\end{itemize}
\subsection{Implementazioni}
Tramite i compomenti elettrici come oscillatori (generano sinusoidi a frequenze desiderate), mixers (moltiplicano i segnali), amplificatori (aumentano la distanza rigenerando il segnale), filtri (selezionano le frequenze) si applicano i concetti logici.

Prima tecnologia tubi a vuoto (delicato, obselescente) surclassata dai transistor (legge di moore\footnote{La legge di moore, del raddopio del numero di transistor ogni 18 mesi, è motivata dal fatto che avere nello stesso spazio più componenti piccole comporta minor costi e maggior resistenza.}).
\subsection{Digitalizzazione}
Le comunicazioni analogiche sono difficili (distorsioni, attenuazioni, interferenze) e richiedono soluzioni specifiche a problemi comuni in base all'implementazione. Soluzione:\[ADC\longleftrightarrow DAC\]
\subsubsection{Fasi digitalizzazione}
Dato in input in segnale analogico le fasi di digitalizzazione sono le seguenti in ordine.
\paragraph{Campionamento} misurazione in predeterminati istanti di tempo che produce una serie temporale di misurazioni.

Possibile operazione distruttiva (impossibile ricostruire segnale originario) con rischio di \textbf{sub-sampling} (aliasing, assenza informazioni) o \textbf{Over-sampling} (valori sovrabbondanti).

Il teorema di Shannon-Nyquist definisce il minimo intervallo che permette di ricostruire il segnale analogico integralmente \[T_s=\frac{1}{f_s}<\frac{1}{2f_M}, f_s>2f_M\]
\paragraph{Quantizzazione} adatta le misurazioni del campionamento al n umero di cifre prederminato per la rappresentazione binaria.

Più bit uso più sarò preciso a rappresentare il segnale, ma sarà sempre "scalinato".
\subsection{Vantaggi digitale}
\paragraph{Integrazione} data dal formato unico dell'informazione, i bit.
\paragraph{Computazione} dei bit naturale per i calcolatori elettronici. Sono facili operazioni di modifica (aggiunta, estrazione, combinazione), compressione e cifratura.
\subsection{Servizio}
La rete permette un certo tipo di comunicazione in base alla tipologia d'interazione nella comunicazione, alla modalità di fluizione dell'informazione e alla tipologia d'informazione.
\begin{itemize}
	\item \textbf{Interattivi}, sorgente e destinazione hanno interazioni. Tipicamente ha uno scambio asincrono dei dati tramite memorizzazione.
	\item \textbf{Distribuiti}, sorgente diffonde informazioni indipendentemente dai destinatari sconosciuti.

	\textbf{Senza controllo di presentazione}, destinatario non controlla ordine di ricezione.

	\textbf{Con controllo di presentazione}, destinatario controlla ordine di ricezione.
\end{itemize}
Il flusso infromativo può essere: \textbf{Punto-Punto (1-1), Multicast (1-N noti), Broadcast (1-tutti), Monodirezionale, Bidirezionale Simmetrico o Asimmetrico}.
Il servizio può essere:
\begin{itemize}
	\item \textbf{Monomediale}, trasferisce un solo tipo d'informazione, oppure più tipologie d'informazione sotto forma d'un solo segnale.
	\item \textbf{Multimediali}, trasferisce più tipi d'informazione e con modalità distinte (per protocolli e qualità di servizio). Da problemi di sincronizzazione.
\end{itemize}
\subsubsection{Qualità del servizio}
\paragraph{La trasparenza} è la capacità della rete di non modificare l'informazione pertinente ad un dialogo.
\begin{itemize}
	\item \textbf{Semantica}, tutela l'integrità dell'informazione trasportata. Gestisce anche situazioni d'errore ripristinando i dati.
	\item \textbf{Temporale}, tutela dai ritardi di transito. Il ritardo è sempre presente per via della propagazione.
\end{itemize}

\textbf{Quality of service (QoS)} è la qualità della comunicazione percepita (funzione della trasparenza).

I sistemi non real-time richiedono alta trasparenza semantica mentre i sistemi real-time trasparenza temporale (servizi isocroni).
\paragraph{La Compensazione al terminale} è una soluzione hai ritardi introdotti dalla rete è applicare un ritardo al terminale del ricevitore che permette di raccogliere in un playback buffer i dati generati dal mittente per ripordurli successivamente dopo un ritardi di playback senza interruzioni.
Si ottiene cosi una trasmissione correttamente temporizzata.

\medskip
\begin{adjustbox}{max width=\textwidth}
	\begin{tikzcd}
		\text{mitt. trasmette} \arrow[rrr, "\text{ritardo rete}" description] &  &  & \text{dest. riceve} \arrow[rrr, "\text{ritardo playback}" description] &  &  & \text{dest. riproduce}
	\end{tikzcd}
\end{adjustbox}

\medskip
Se il tempo ri ricezione coincide con il tempo di riproduzione si verifica un interruzione del servizio. (pensa a video in streaming con barra rossa e barra grigia).
\subparagraph{Le reti CBR e VBR} son definire a flusso costante e variabile di bit.
\subparagraph{Le reti a tempo differito e non} sono quando il dato viene usato qunado il destinatario vuole e qunado va consumato subito (trasmissioni in line).
\subsubsection{Multiplazione}
La multiplazione è il trasporto di più canali sullo stesso mezzo trasmissivo.

Multiplazione=una rete per più servizi con requisiti diversi. Serve una rete flessibile in allocazione, distribuisce le risorse secondo le necessità, e gestione, richieste di servizi diversi non interferiscono tra loro.

\paragraph{Le diverse realizzazioni} sono teoricamente, in condizioni ottimali, equivalenti:
\begin{itemize}
	\item \textbf{FDM} basata sulle frequenze usa un approccio analogico. Hai canali digitali viene assegnato un differente intervallo di frequenze.
	\item \textbf{TDM} basata sul tempo usa un approccio digitale. L'intero mezzo trasmissivo viene usato a giro da un solo canale per una finestra temporale fissa.
	\item \textbf{CDM} basata su codici ortogonali associati hai canali che trasmettono simultaneamente e sulla stessa frequenza.

	Il mittente moltiplica il segnale per il codice del ricevitore e i ricevitori moltiplicheranno il messaggio ricevuto per il proprio codici ortogonali. L'unico ricevitore avrà un risultato diverso da zero è quello a cui era indirizzato il messaggio.
	\item Divisione di spazio.
\end{itemize}
\subsubsection{TDM}
La TDM nelle reti digitali ha il maggior successo.
Casi di applicazioni della TDM con assegnamento della banda staticamente o dinamicamente:

\medskip
\begin{adjustbox}{max width=\textwidth}
	\begin{tikzcd}
		& Slotted \arrow[rd] \arrow[r]                           & framed \arrow[rr, "usata" description] \arrow[rrd, "inusata" description] &  & statica  \\
		\textbf{TDM} \arrow[ru] \arrow[r] & unslotted \arrow[rrr, "usata" description, bend right] & unframed \arrow[rr, "inusata" description]                                &  & dinamica
	\end{tikzcd}
\end{adjustbox}
\paragraph{slotted} il tempo è diviso in slot di dimensione fissa. Le unità informative sono hanno la stessa lunghezza e commisurata allo slot.
	\subparagraph{slotted framed} struttura gli slot in frame di dimensioni prefissate e sincronizzati da un slot dedicato.

Ogni slot in un frame ha una posizione e è un canale a sè. La posizione si ripete tra frame ma in flussi separati (stesso codice canali diversi).
\subparagraph{slotted unframed (S-TDM)} gestisce sequenze di slot di lunghezza variabile tramite slot di sincronizzazione.

Ogni slot in una sequenza è numerato e è un canale a sè. La numerazione si ripete tra sequenze ma in flussi separati (stesso codice canali diversi).
\paragraph{unslotted} delimita le unità informative di lunghezza variabile tramite delimitatori.
\paragraph{A.B.statica} il flusso informativo coincide alla banda  dedicatagli, che non può cambiare durante la comunicazione.
\paragraph{A.B.dinamica (A-TDM)} il flusso informativo condivide la banda con altri flussi, secondo la necessità. La banda può cambiare durante la  comunicazione.

Necessario definire come si assegna la banda e come gestire situazioni di contesa.
\subsection{La rete}
La rete consente la comunicazione tra una combinazione qualunque di elementi terminali, tramite canali end-to-end dinamicamente creati all'occorrenza, con una qualità del servizio accettabile.

\paragraph{End-to-end} è il canale supportato da diversi mezzi trasmissivi, in funzione dell'architettura della rete, per connette sorgente e destinazione.
\paragraph{ipercubo} La rete è costituita dal piano di gestione della rete più Piano di controllo della rete e i dati di controllo. Quest osecondo contiene i dati d'utente.
\subsubsection{Funzionalità}
\begin{multicols}{2}
	\paragraph{Trasmissione} gestisce il trasferimento fisico del segnale punto-punto/punto-multipunto.
	\paragraph{Commutazione} gestisce l'instradamento delle informazioni nella rete per far comunicare i nodi terminali.
	\paragraph{Segnalazione} gestisce lo scambio di informazioni della rete stessa e per gestire la comunicazione.
	\paragraph{Gestione} tutto ciò che riguarda il mantenimento delle funzioni della rete.
\end{multicols}
\subsubsection{Commutazione}
\paragraph{Circuito} gestisce informazioni analogiche e digitali.
La rete crea un cananale dedicato e esclusivo tra i due nodi terminali.

\medskip
\begin{adjustbox}{max width=\textwidth}
	\begin{tikzcd}
		Instaurazione \ del \ circuito \arrow[rr] &  & Dialogo \arrow[rr] &  & Disconnessione \ del \ circuito
	\end{tikzcd}
\end{adjustbox}

\subparagraph{Call set-up time} è un ritardo iniziale per settare il canale.
Rete trasparente temporalmente durante il dialogo.
Necessaria disconnessione al circuito per liberare le risorse di rete.
\paragraph{Messaggio} gestisce informazioni digitali.
Struttura le informazioni d'utente e di segnalazione in messaggi suddivisi in pacchetti di dimensione prefissata.
\subparagraph{Circuito virtuale} prestabilito percorso dei pacchetti a cui è associato un numero di circuito virtuale. I pacchetti hanno solo il numero di circuito.
\subsection{Datagram} pacchetti instradati indipendentemente con informazioni d'indirizzamento. Possono esserci tempo d'arrivo differenti che comportano un riordinamento dei pacchetti a destinazione.